% ============================================================
%  MDGen: Generative Molecular Dynamics Simulation
%  A Comprehensive Technical Survey
%  Author: Ditsa Mallick  ·  dmallic1@asu.edu
% ============================================================
\documentclass[11pt,a4paper]{article}

% --- Packages -----------------------------------------------
\usepackage[T1]{fontenc}
\usepackage[utf8]{inputenc}
\usepackage{lmodern}
\usepackage[margin=1in]{geometry}
\usepackage{amsmath,amssymb,amsthm}
\usepackage{bm}
\usepackage{booktabs}
\usepackage{array}
\usepackage{multirow}
\usepackage{tabularx}
\usepackage{graphicx}
\usepackage{xcolor}
\usepackage{hyperref}
\usepackage{microtype}
\usepackage{setspace}
\usepackage{parskip}
\usepackage{titlesec}
\usepackage{fancyhdr}
\usepackage{algorithm}
\usepackage{algpseudocode}
\usepackage{enumitem}
\usepackage{caption}
\usepackage{subcaption}
\usepackage{cite}
\usepackage{url}
\usepackage{mathtools}
\usepackage{tcolorbox}
\usepackage{mdframed}
\usepackage{wrapfig}
\usepackage{float}

% --- Colour palette -----------------------------------------
\definecolor{teal}{RGB}{11,114,133}
\definecolor{rust}{RGB}{194,65,12}
\definecolor{amber}{RGB}{240,140,0}
\definecolor{ink}{RGB}{16,37,34}
\definecolor{muted}{RGB}{83,98,94}
\definecolor{bglight}{RGB}{232,243,239}

% --- Hyperref setup -----------------------------------------
\hypersetup{
  colorlinks=true,
  linkcolor=teal,
  citecolor=teal,
  urlcolor=rust,
  pdftitle={MDGen: Generative Molecular Dynamics Simulation -- A Comprehensive Survey},
  pdfauthor={Ditsa Mallick},
}

% --- Page style ---------------------------------------------
\pagestyle{fancy}
\fancyhf{}
\lhead{\small\textcolor{muted}{MDGen: Generative Molecular Dynamics Simulation}}
\rhead{\small\textcolor{muted}{\thepage}}
\renewcommand{\headrulewidth}{0.4pt}

% --- Section formatting -------------------------------------
\titleformat{\section}{\large\bfseries\color{ink}}{{\thesection}}{0.8em}{}[\color{teal}\rule{\linewidth}{0.5pt}]
\titleformat{\subsection}{\normalsize\bfseries\color{ink}}{{\thesubsection}}{0.8em}{}
\titleformat{\subsubsection}{\normalsize\itshape\color{muted}}{{\thesubsubsection}}{0.8em}{}

% --- Custom environments ------------------------------------
\newmdenv[linecolor=teal,backgroundcolor=bglight!50,
  leftmargin=0pt,rightmargin=0pt,
  skipabove=8pt,skipbelow=8pt,
  innerleftmargin=10pt,innerrightmargin=10pt,
  innertopmargin=6pt,innerbottommargin=6pt]{calloutbox}

\newmdenv[linecolor=rust,backgroundcolor=rust!5,
  leftmargin=0pt,rightmargin=0pt,
  skipabove=8pt,skipbelow=8pt,
  innerleftmargin=10pt,innerrightmargin=10pt,
  innertopmargin=6pt,innerbottommargin=6pt]{warningbox}

% --- Math operators -----------------------------------------
\DeclareMathOperator{\SE}{SE}
\DeclareMathOperator{\SO}{SO}
\DeclareMathOperator{\IGSO}{IGSO3}
\newcommand{\bx}{\mathbf{x}}
\newcommand{\bX}{\mathbf{X}}
\newcommand{\bs}{\mathbf{s}}
\newcommand{\br}{\mathbf{r}}
\newcommand{\bR}{\mathbf{R}}
\newcommand{\bt}{\mathbf{t}}
\newcommand{\bW}{\mathbf{W}}
\newcommand{\beps}{\boldsymbol{\epsilon}}
\newcommand{\RR}{\mathbb{R}}
\newcommand{\TT}{\mathbb{T}}
\newcommand{\EE}{\mathbb{E}}
\newcommand{\NN}{\mathcal{N}}

% ============================================================
\begin{document}
% ============================================================

% --- Title page ---------------------------------------------
\begin{titlepage}
  \centering
  \vspace*{2cm}

  {\Large\bfseries\color{teal} Technical Survey}\\[0.4em]
  \rule{\textwidth}{1pt}\\[0.6em]

  {\LARGE\bfseries\color{ink}
    MDGen: Generative Molecular Dynamics Simulation\\
    \large A Comprehensive Technical Survey with Background, Prior Work, and Analysis}\\[0.6em]

  \rule{\textwidth}{1pt}\\[1.5cm]

  {\large Ditsa Mallick}\\[0.3em]
  {\normalsize School of Computing and Augmented Intelligence\\
   Arizona State University, Tempe, AZ 85281}\\[0.2em]
  {\normalsize \texttt{dmallic1@asu.edu}}\\[1.5cm]

  {\normalsize
    \textbf{Based on:} Jing, B., St\"{a}rk, H., Jaakkola, T., \& Berger, B. (2024).\\
    Generative Molecular Dynamics. \textit{arXiv:2409.17808}. NeurIPS 2024.}\\[1cm]

  {\normalsize May 2026}\\[2cm]

  \begin{abstract}
    \noindent
    Molecular dynamics (MD) simulation is the gold standard for studying protein conformational
    dynamics, yet its computational cost---requiring hundreds of millions of femtosecond integration
    steps to probe biologically relevant timescales---severely limits its applicability to large-scale
    protein screening and engineering campaigns.  This survey provides a comprehensive technical
    exposition of \textbf{MDGen} (Jing et al., NeurIPS 2024), the first generative model to produce
    \emph{temporally coherent}, full-backbone protein trajectories at approximately
    $10{,}000\times$ the speed of equivalent explicit-solvent MD.  We begin with the necessary
    background in protein biophysics, statistical mechanics, and the theory of denoising diffusion
    probabilistic models, including their extension to non-Euclidean manifolds such as $\SO(3)$ and
    the torus $\TT^4$.  We then survey the landscape of prior computational approaches---Normal Mode
    Analysis, Boltzmann generators, independent-frame diffusion models, and equilibrium ensemble
    methods---and characterise their shared failure to capture temporal kinetics.  The MDGen
    architecture, which couples SE(3)-equivariant Invariant Point Attention across residues with
    axial temporal attention across trajectory frames, is described in full mathematical detail.
    Training on the ATLAS database of 1,400 diverse 300~ns protein simulations and evaluation on
    held-out test proteins demonstrate RMSF Pearson correlation $r \approx 0.87$ (versus Normal
    Mode Analysis $r \approx 0.64$), physically correct Ramachandran distributions
    (Jensen--Shannon divergence $\approx 0.08$), and correct frame-to-frame step-size and
    autocorrelation function profiles---metrics that independent samplers cannot reproduce.
    We close with discussion of applications to cryptic pocket discovery and intrinsically disordered
    proteins, known limitations, and promising future directions.
  \end{abstract}
\end{titlepage}

% --- Table of contents --------------------------------------
\tableofcontents
\newpage

% ============================================================
\section{Introduction}
% ============================================================

Proteins are the primary molecular machines of living systems.  They catalyse chemical reactions,
transmit signals, provide structural support, and execute the vast majority of the biochemical
functions required for life.  A protein's biological activity depends critically not only on its
average three-dimensional structure but on its \emph{dynamical} behaviour: the ensemble of
conformations it explores at physiological temperature, the timescales of interconversion between
states, and the kinetic pathways connecting them~\cite{adcock2006}.

Despite decades of progress in experimental structure determination---X-ray crystallography,
cryo-electron microscopy, NMR spectroscopy---these methods report static or ensemble-averaged
snapshots.  Molecular dynamics (MD) simulation fills this gap by numerically integrating Newton's
equations of motion for every atom, generating detailed time-series of atomic positions.  Modern
GPU-accelerated MD packages (AMBER~\cite{case2021}, GROMACS~\cite{abraham2015}, OpenMM) can
simulate a 100-residue protein at roughly 100--200~ns per day.  Yet the biologically important
timescales---domain motions, loop gating, and folding transitions---occur on microsecond to
millisecond timescales, spanning 9--12 orders of magnitude above the obligatory 2~fs integration
time step imposed by the fastest atomic vibrations.  This \emph{timescale gap} makes MD
prohibitively expensive for large-scale ensemble characterisation, mutational screening, or drug
discovery campaigns requiring thousands of protein variants.

\begin{calloutbox}
  \textbf{Core motivation.}  The central question addressed by MDGen is: \emph{can a neural
  network, trained on a large corpus of MD trajectories, learn to generate new trajectory segments
  that faithfully reproduce both the equilibrium thermodynamics and the kinetic temporal structure
  of the original simulations---without running the simulation itself?}
\end{calloutbox}

The emergence of deep generative models---particularly denoising diffusion probabilistic
models~\cite{ho2020} and score-based continuous-time formulations~\cite{song2021}---has
transformed structure-prediction~\cite{jumper2021}, protein design~\cite{watson2023}, and
equilibrium ensemble generation~\cite{jing2023eigenfold,jing2024alphaflow}.  However, all prior
generative models for protein conformations operate in an \emph{independent, identically
distributed} (i.i.d.)\ sampling regime: each sample is drawn independently from an approximation
of the equilibrium distribution.  Such models cannot, by construction, reproduce temporal
autocorrelation functions, kinetically realistic step sizes between consecutive frames, or any
other observable that depends on the \emph{ordering} of configurations in time.

MDGen~\cite{jing2024mdgen} addresses this gap by posing the generative problem at the level of
entire trajectory \emph{segments}: given a protein sequence, generate a window of $T$ consecutive
backbone frames whose joint distribution matches that of a real MD simulation.  The key technical
contribution is an architecture that alternates SE(3)-equivariant \emph{spatial} attention---the
Invariant Point Attention (IPA) module from AlphaFold2~\cite{jumper2021}---with a novel
\emph{temporal axial attention} layer that propagates information between frames along the time
axis.  Trained on the ATLAS database~\cite{vandermeersche2024} of 1,400 diverse proteins with
300~ns of explicit-solvent MD each, MDGen generates trajectories at approximately $10{,}000\times$
the speed of equivalent MD while achieving a RMSF Pearson correlation of $r \approx 0.87$ and
reproducing correct temporal kinetics.

This survey provides a self-contained technical exposition of MDGen accessible to readers with a
machine learning background but limited exposure to structural biology, and vice versa.  We
intentionally proceed from first principles.  Section~\ref{sec:background} introduces protein
structure, molecular dynamics, and statistical mechanics.  Section~\ref{sec:diffusion} covers the
theory of diffusion models and their extension to non-Euclidean geometries.
Section~\ref{sec:priorwork} surveys prior computational approaches and their limitations.
Section~\ref{sec:mdgen} presents the MDGen method in full detail.  Sections~\ref{sec:data} and
\ref{sec:results} describe the training data and experimental results.  Section~\ref{sec:applications}
discusses applications, Section~\ref{sec:limitations} limitations, and Section~\ref{sec:future}
future directions.


% ============================================================
\section{Biological and Physical Background}
\label{sec:background}
% ============================================================

\subsection{Proteins: Structure and Function}

A protein is a linear polymer of \emph{amino acid residues} connected by peptide bonds.  Twenty
standard amino acids, differing only in their \emph{side chain} (the group branching from the
$\alpha$-carbon, \smash{C$_\alpha$}), are encoded by the genetic code.  Proteins range from tens
to thousands of residues in length; a typical globular protein has 100--500 residues.

The structural hierarchy of proteins is conventionally decomposed into four levels:
\begin{itemize}[leftmargin=2em]
  \item \textbf{Primary structure:} the amino-acid sequence, a string over a 20-letter alphabet.
  \item \textbf{Secondary structure:} local regularities of the backbone---$\alpha$-helices
    (right-handed coils stabilised by $i \to i+4$ hydrogen bonds) and $\beta$-strands (extended
    chains forming hydrogen-bonded $\beta$-sheets).
  \item \textbf{Tertiary structure:} the complete three-dimensional fold of a single polypeptide
    chain, governed by a combination of hydrophobic packing, electrostatics, hydrogen bonding, and
    van der Waals interactions.
  \item \textbf{Quaternary structure:} the assembly of multiple polypeptide chains into a complex.
\end{itemize}

The protein \textbf{backbone} consists of the repeating $\text{N}$--$\text{C}_\alpha$--$\text{C}{=}\text{O}$
unit that runs through every residue.  The backbone conformation is almost completely determined by
two dihedral angles per residue: $\phi$ (rotation about the N--$\text{C}_\alpha$ bond) and $\psi$
(rotation about the $\text{C}_\alpha$--C bond).  MDGen operates exclusively on backbone heavy atoms
(N, $\text{C}_\alpha$, C, O), representing each residue as a rigid frame, with side chains encoded
separately as torsion angles.

\subsection{The Protein Energy Landscape and Conformational Ensembles}

A protein in solution at temperature $T$ does not adopt a single static conformation.  Instead, it
continuously fluctuates among the set of energetically accessible states according to Boltzmann
statistics.  The probability of finding the system in a conformation $\bx$ (specifying all atomic
coordinates) is given by the canonical ensemble distribution:
\begin{equation}
  p(\bx) \propto \exp\!\left(-\frac{E(\bx)}{k_B T}\right),
  \label{eq:boltzmann}
\end{equation}
where $E(\bx)$ is the potential energy of the conformation, $k_B$ is Boltzmann's constant, and
$k_B T \approx 0.6$~kcal/mol at physiological temperature (300~K).

The \emph{energy landscape} of a protein~\cite{frauenfelder1991} is a high-dimensional surface
in conformation space.  Stable conformations correspond to energy basins (local minima).
Transitions between basins require thermal fluctuations to overcome energy barriers; the rate of
crossing a barrier of height $\Delta E^\ddagger$ scales as $k \propto \exp(-\Delta E^\ddagger /
k_BT)$ by the Arrhenius equation.  Key features of the energy landscape relevant to MDGen include:
\begin{enumerate}[leftmargin=2em]
  \item \textbf{Multiple minima:} most functional proteins populate more than one conformation
    (e.g.\ open/closed states of enzyme active sites, folded/partially unfolded states).
  \item \textbf{Hierarchical roughness:} fast local motions (bond vibrations, 10~fs) are nested
    within slower large-scale rearrangements ($\mu$s--ms).
  \item \textbf{Anharmonicity:} the landscape is decidedly non-quadratic; large-amplitude motions
    cannot be described by a harmonic approximation.
\end{enumerate}

The biologically relevant object is thus the \emph{conformational ensemble}---the full distribution
$p(\bx)$---not a single representative structure.  Dynamical properties such as the rate of
conformational change, the timescales of domain motion, and the population of cryptic binding
pockets all depend on the shape of this distribution and the barriers between its modes.

\subsection{Molecular Dynamics Simulation}

Molecular dynamics (MD) simulation~\cite{adcock2006} generates samples from the conformational
ensemble $p(\bx)$ by numerically integrating Newton's equations of motion:
\begin{equation}
  m_i \frac{d^2 \br_i}{dt^2} = \bm{F}_i = -\nabla_{\br_i} E(\br_1, \ldots, \br_N),
  \label{eq:newton}
\end{equation}
where $m_i$ and $\br_i$ are the mass and position of atom $i$, and $E$ is the empirical
\emph{force field} potential energy function.

The force field $E$ comprises:
\begin{itemize}[leftmargin=2em]
  \item \textbf{Bonded terms:} harmonic bond stretches, angle bends, and periodic dihedral torsions
    calibrated to quantum chemical calculations.
  \item \textbf{Non-bonded terms:} the Lennard-Jones 6-12 potential for van der Waals interactions,
    and Coulomb electrostatics (often treated via Particle Mesh Ewald for long-range contributions).
\end{itemize}
Widely used force fields include AMBER ff14SB~\cite{maier2015}---the force field used for the ATLAS
dataset on which MDGen is trained---CHARMM36m, and OPLS-AA.

\textbf{The timescale problem.}  The fastest atomic vibrations (bond stretches, $\sim$10~fs) impose
an integration time step ceiling of $\Delta t = 2$~fs (with bond-length constraints).  Observing
biologically important motions requires $10^8$--$10^{11}$ steps, corresponding to
100~ns--100~$\mu$s of simulation time.  A single microsecond simulation of a 100-residue protein
in explicit solvent on a modern A100 GPU cluster requires roughly 1--2 weeks of wall-clock time.
Screening thousands of protein variants---as required in drug discovery or protein engineering---is
therefore practically infeasible with standard MD.

\textbf{Enhanced sampling methods.}  Multiple strategies have been developed to alleviate the
timescale problem: replica exchange MD~\cite{sugita1999} runs parallel copies at different
temperatures; metadynamics~\cite{laio2002} applies a history-dependent bias potential to accelerate
barrier crossings; accelerated MD~\cite{hamelberg2004} lowers barriers globally.  These methods
improve sampling of rare events but still require explicit simulation and do not provide the
generative, sequence-conditioned model that MDGen offers.

\subsection{Trajectory Representation}

An MD trajectory is a time-ordered sequence of \emph{frames} $(\bx_1, \bx_2, \ldots, \bx_T)$,
where each frame $\bx_t$ specifies the full set of atomic coordinates at time $t$.  In MDGen, the
relevant degrees of freedom per residue $i$ at time $t$ are:
\begin{itemize}[leftmargin=2em]
  \item $(\bR_i^{(t)}, \bt_i^{(t)}) \in \SE(3)$: a rotation (backbone orientation) and
    translation ($\text{C}_\alpha$ position) defining a rigid body frame.
  \item $\bm\chi_i^{(t)} \in \TT^4$: the four side-chain torsion angles on the 4-torus.
\end{itemize}
The full state of a protein with $N$ residues at one time step is:
\begin{equation}
  \bx = \{(\bR_i, \bt_i, \bm\chi_i)\}_{i=1}^N \in \SE(3)^N \times \TT^{4N},
\end{equation}
and a trajectory window of $T$ frames is:
\begin{equation}
  \bX = (\bx_1, \bx_2, \ldots, \bx_T) \in \bigl(\SE(3)^N \times \TT^{4N}\bigr)^T.
  \label{eq:trajectory}
\end{equation}


% ============================================================
\section{Theoretical Background: Diffusion Models}
\label{sec:diffusion}
% ============================================================

\subsection{Denoising Diffusion Probabilistic Models (DDPMs)}

Denoising diffusion probabilistic models (DDPMs)~\cite{ho2020} define a generative model through
two coupled Markov chains: a \emph{forward process} that gradually destroys structure in the data
by adding Gaussian noise, and a \emph{reverse process} that is learned to reconstruct data from
noise.

\textbf{Forward process.}  Given a clean data sample $\bx_0 \sim q_{\text{data}}(\bx)$, the
forward process $q(\bx_{1:S} \mid \bx_0)$ adds Gaussian noise over $S$ discrete steps:
\begin{equation}
  q(\bx_s \mid \bx_{s-1}) = \mathcal{N}\!\left(\bx_s;\; \sqrt{1-\beta_s}\,\bx_{s-1},\;
  \beta_s \bm{I}\right),
  \label{eq:forward_step}
\end{equation}
where $\beta_s \in (0, 1)$ is a pre-specified noise schedule.  Crucially, the marginal at step $s$
has a closed form:
\begin{equation}
  q(\bx_s \mid \bx_0) = \mathcal{N}\!\left(\bx_s;\; \sqrt{\bar\alpha_s}\,\bx_0,\;
  (1-\bar\alpha_s)\bm{I}\right),
  \label{eq:forward_marginal}
\end{equation}
where $\bar\alpha_s = \prod_{k=1}^{s} (1-\beta_k)$.  As $s \to S$, $\bar\alpha_s \to 0$ and the
distribution approaches an isotropic Gaussian $\mathcal{N}(\bm{0}, \bm{I})$.

\textbf{Reverse process.}  The reverse process $p_\theta(\bx_{s-1} \mid \bx_s)$ is parameterised
as a Gaussian:
\begin{equation}
  p_\theta(\bx_{s-1} \mid \bx_s) = \mathcal{N}\!\left(\bx_{s-1};\;
  \bm\mu_\theta(\bx_s, s),\; \sigma_s^2 \bm{I}\right).
\end{equation}
Ho et al.~\cite{ho2020} showed that $\bm\mu_\theta$ can be re-parameterised in terms of a
\emph{noise-prediction network} $\beps_\theta$:
\begin{equation}
  \bm\mu_\theta(\bx_s, s) = \frac{1}{\sqrt{\alpha_s}}\!\left(\bx_s -
  \frac{\beta_s}{\sqrt{1-\bar\alpha_s}}\,\beps_\theta(\bx_s, s)\right),
\end{equation}
and the training objective simplifies to:
\begin{equation}
  \mathcal{L}_{\text{DDPM}} = \EE_{s, \bx_0, \beps}\!\left[
  \|\beps - \beps_\theta(\sqrt{\bar\alpha_s}\,\bx_0 + \sqrt{1-\bar\alpha_s}\,\beps,\; s)\|^2
  \right],
  \label{eq:ddpm_loss}
\end{equation}
where $\beps \sim \mathcal{N}(\bm{0}, \bm{I})$ is the noise added to $\bx_0$.  Minimising this
objective is equivalent to learning to predict which noise was added, which in turn corresponds to
learning the \emph{score function} $\nabla_{\bx_s} \log p_s(\bx_s)$ scaled by a factor that
depends on the noise schedule.

\subsection{Score-Based Generative Models and the Continuous-Time SDE Perspective}

Song et al.~\cite{song2021} cast diffusion models as continuous-time stochastic differential
equations (SDEs), providing a unified theoretical framework.  The forward SDE is:
\begin{equation}
  d\bx = f(\bx, t)\,dt + g(t)\,d\bW,
  \label{eq:forward_sde}
\end{equation}
where $\bW$ is a standard Brownian motion, $f$ is a drift coefficient, and $g$ is a diffusion
coefficient.  The reverse SDE~\cite{anderson1982}, which runs backwards in time from noise to data,
is:
\begin{equation}
  d\bx = \left[f(\bx, t) - g(t)^2\,\nabla_\bx \log p_t(\bx)\right]dt + g(t)\,d\bar\bW,
  \label{eq:reverse_sde}
\end{equation}
where $\bar\bW$ is a reverse-time Brownian motion.  The key quantity is the
\textbf{score function} $\nabla_\bx \log p_t(\bx)$, which points toward regions of higher
probability density.  A \emph{score network} $s_\theta(\bx, t)$ is trained via
\textbf{denoising score matching}~\cite{vincent2011} to approximate this score:
\begin{equation}
  \mathcal{L}_{\text{DSM}} = \EE_{t, \bx_0, \bx_t}\!\left[\omega(t)\,
  \|s_\theta(\bx_t, t) - \nabla_{\bx_t}\log q(\bx_t \mid \bx_0)\|^2\right],
\end{equation}
where $\omega(t)$ is a weighting function and the expectation is taken over the forward process.
The optimal $s_\theta$ equals $\nabla_\bx \log p_t(\bx_t)$, the true score.

\textbf{DDIM sampling.}  Denoising Diffusion Implicit Models (DDIM)~\cite{song2021ddim} provide a
deterministic ODE equivalent of the reverse SDE, enabling high-quality sample generation with far
fewer function evaluations (typically 50 steps instead of 500):
\begin{equation}
  \bx_{s-1} = \sqrt{\bar\alpha_{s-1}}\!\left(\frac{\bx_s - \sqrt{1-\bar\alpha_s}\,\beps_\theta}
  {\sqrt{\bar\alpha_s}}\right) + \sqrt{1-\bar\alpha_{s-1}}\,\beps_\theta.
  \label{eq:ddim}
\end{equation}

\subsection{Diffusion on Non-Euclidean Manifolds}

Standard Gaussian diffusion assumes data lies in flat Euclidean space $\RR^d$.  Protein backbone
frames involve two non-Euclidean components:
\begin{enumerate}[leftmargin=2em]
  \item \textbf{Rotations} $\bR_i \in \SO(3)$, a compact 3D Lie group with non-trivial curved
    geometry.
  \item \textbf{Torsion angles} $\chi_j \in S^1$ (the circle), which wrap around with period
    $2\pi$; the side-chain state lives on the 4-torus $\TT^4 = (S^1)^4$.
\end{enumerate}

\textbf{Diffusion on $\SO(3)$: the IGSO3 distribution.}  De Bortoli et al.~\cite{debortoli2022}
and Leach et al.~\cite{leach2022} developed Riemannian diffusion models for rotation groups.
MDGen uses the \textbf{Isotropic Gaussian on $\SO(3)$} (IGSO3):
\begin{equation}
  q(\bR_s \mid \bR_0) = \IGSO(\ \bR_s;\; \bR_0,\; \sigma_s^2\ ),
  \label{eq:igso3}
\end{equation}
whose score involves the Riemannian gradient---the analogue of $\nabla_\bx \log p$ on the curved
$\SO(3)$ manifold.  Concretely, if $\bR_s = \exp_{\bR_0}(\bm\xi)\,\bR_0$ where $\bm\xi \in
\mathfrak{so}(3)$ is in the Lie algebra (the tangent space at the identity), then the score in the
tangent space at $\bR_s$ can be expressed via the heat kernel on $\SO(3)$~\cite{debortoli2022}.

\textbf{Diffusion on the torus $\TT$: wrapped Gaussians.}  For torsion angles $\theta \in
(-\pi, \pi]$, the correct diffusion distribution is a \emph{wrapped Gaussian}~\cite{jing2022torsional}:
\begin{equation}
  q(\theta_s \mid \theta_0) = \sum_{k=-\infty}^{\infty} \mathcal{N}(\theta_s;\;
  \theta_0 + 2\pi k,\; \sigma_s^2),
  \label{eq:wrapped_gaussian}
\end{equation}
which accounts for the periodicity of the angular domain.  Denoising score matching on $\TT$ is
computed using the wrapped score---the derivative of the log wrapped Gaussian density with respect
to $\theta_s$, which can be evaluated by truncating the infinite sum.

\subsection{SE(3) Equivariance}

A fundamental symmetry of physical systems in three-dimensional space is \emph{Euclidean
equivariance}.  The group $\SE(3) = \SO(3) \ltimes \RR^3$ consists of all rotations and
translations of $\RR^3$.  A map $f$ is \emph{SE(3)-equivariant} if:
\begin{equation}
  f(g \cdot \bx) = \rho_{\text{out}}(g)\,f(\bx), \quad g \in \SE(3),
  \label{eq:equivariance}
\end{equation}
where $\rho_{\text{out}}(g)$ is a representation of $g$ acting on the output space.  For
coordinate outputs, this means that rotating the input configuration rotates the output coordinates
by the same rotation.  For scalar outputs (e.g.\ energy), SE(3)-equivariance reduces to
\emph{invariance}: $f(g \cdot \bx) = f(\bx)$.

SE(3)-equivariant neural networks~\cite{fuchs2020,satorras2021} are architecturally constrained so
that~\eqref{eq:equivariance} holds exactly.  This is essential for molecular modelling: the physics
of a protein does not change if we rotate or translate the entire coordinate system.  A model that
violates this symmetry must learn it from data, which wastes capacity and may fail to generalise.

\subsection{Backbone Frame Representation}

The AlphaFold2 structure module~\cite{jumper2021} introduced the \emph{backbone frame}
representation that MDGen inherits.  Each residue $i$ is assigned a rigid body in $\SE(3)$:
\begin{itemize}[leftmargin=2em]
  \item \textbf{Translation} $\bt_i \in \RR^3$: the position of the $\text{C}_\alpha$ atom.
  \item \textbf{Rotation} $\bR_i \in \SO(3)$: a local coordinate frame defined by the N,
    $\text{C}_\alpha$, and C atoms of the residue backbone, encoding backbone orientation.
\end{itemize}
Together, $T_i = (\bR_i, \bt_i)$ defines a frame in which the local chemical environment of
residue $i$ appears in a standardised orientation.  The full protein backbone is then represented
as a set of $N$ frames $\{T_i\}_{i=1}^N \in \SE(3)^N$.


% ============================================================
\section{Prior Work and Its Limitations}
\label{sec:priorwork}
% ============================================================

Understanding the novelty of MDGen requires surveying the space of prior methods for computational
characterisation of protein conformational dynamics.  We organise this survey into four categories:
classical physics-based methods, equilibrium ML samplers, single-structure ML predictors, and
methods that partially address ensemble generation.

\subsection{Normal Mode Analysis}

Normal Mode Analysis (NMA)~\cite{tirion1996,brooks1983} approximates the protein energy landscape
as a multidimensional quadratic (harmonic) bowl centred on a known minimised structure.  Under this
approximation, the equations of motion separate into $3N-6$ independent harmonic oscillators (normal
modes), each with a characteristic frequency.  The lowest-frequency modes correspond to the
largest-scale, most collective motions and can be used to predict B-factors and RMSF profiles.

\textbf{Strengths.}
\begin{itemize}[leftmargin=2em]
  \item Computationally cheap: a single eigendecomposition of the Hessian matrix.
  \item Provides physically interpretable modes (e.g.\ a ``breathing'' mode of a viral capsid).
  \item Reasonable agreement with MD-derived RMSF for well-ordered secondary structure.
\end{itemize}

\textbf{Fundamental limitations.}
\begin{itemize}[leftmargin=2em]
  \item \textbf{Harmonic approximation:} NMA is strictly valid only near the energy minimum.
    Loop motions, domain swings, and any transition requiring barrier crossing are
    \emph{outside the scope of the harmonic model}.
  \item \textbf{No kinetics:} NMA produces independent samples from a Gaussian approximation to
    the equilibrium ensemble; it has no notion of temporal ordering or transition rates.
  \item \textbf{Single minimum:} NMA cannot represent multi-modal distributions or conformational
    heterogeneity between distinct basins.
  \item In direct comparison on the ATLAS test set, NMA achieves RMSF Pearson correlation
    $r \approx 0.64$---substantially below MDGen's $r \approx 0.87$.
\end{itemize}

\subsection{Boltzmann Generators and Flow-Based Equilibrium Samplers}

Boltzmann generators~\cite{noe2019} use normalising flows (bijective, differentiable maps) to
learn a direct mapping between a tractable base distribution (Gaussian) and the Boltzmann
distribution $p(\bx) \propto \exp(-E(\bx)/k_BT)$.  At test time, samples can be drawn from the
Boltzmann distribution without simulating the dynamics.

\textbf{Strengths.}
\begin{itemize}[leftmargin=2em]
  \item Principled probabilistic model with exact likelihood.
  \item Can, in principle, provide reweighting factors to correct for force-field imperfections.
  \item Demonstrated on small peptides and miniproteins.
\end{itemize}

\textbf{Fundamental limitations.}
\begin{itemize}[leftmargin=2em]
  \item \textbf{Scalability:} normalising flows for all-atom protein configurations of 100--500
    residues remain computationally intractable; most demonstrations are restricted to
    $\leq$50-residue systems.
  \item \textbf{No temporal coherence:} like NMA, Boltzmann generators produce i.i.d.\ equilibrium
    samples.  There is no mechanism to generate temporally ordered trajectories.
  \item \textbf{Requires the energy function:} the training procedure typically requires
    evaluating $E(\bx)$, which couples the method to the expensive force field.
\end{itemize}

\subsection{Independent-Frame Protein Diffusion Models}

A family of recent diffusion models generates protein structures by learning to reverse a noising
process on backbone frames.

\textbf{FrameDiff}~\cite{yim2023framediff} diffuses independently over the $\SE(3)^N$ space of
backbone frames and trains a score network that denoises them.  At inference, T samples are
drawn \emph{independently} to form a structural ensemble.

\textbf{EigenFold}~\cite{jing2023eigenfold} decomposes the protein structure along a hierarchy of
eigenmodes of a spring network, performing diffusion along each eigenmode independently.  This
provides a natural coarse-to-fine generation scheme.

\textbf{AlphaFlow}~\cite{jing2024alphaflow} fine-tunes the AlphaFold2 structure module as a
flow-matching model conditioned on multiple-sequence alignments (MSAs), enabling
ensemble prediction for proteins with known evolutionary context.

\textbf{BioEmu}~\cite{lewis2025bioemu} trains a large-scale diffusion model on both MD trajectories
and experimental data, yielding strong equilibrium ensemble predictions.

\textbf{Shared limitation: temporal incoherence.}  All of the above methods generate \emph{independent}
conformational samples.  While they may approximate the marginal equilibrium distribution $p(\bx)$,
they cannot reproduce any temporal property:
\begin{enumerate}[leftmargin=2em]
  \item The \textbf{autocorrelation function} $C(\tau) = \langle \bx(t) \cdot \bx(t+\tau) \rangle /
    \langle \bx(t)^2 \rangle$ drops to zero immediately at lag $\tau \geq 1$ for i.i.d.\
    samplers, whereas physical trajectories decay smoothly over many nanoseconds.
  \item \textbf{Frame-to-frame step size}: consecutive frames in real MD differ by only
    $\sim$0.1--0.3~\AA\ RMSD.  Independent samples show arbitrarily large jumps with no
    physical regularisation.
  \item \textbf{Transition pathway information}: the sequence of conformations visited during a
    domain opening, the order of secondary structure formation---all are invisible to i.i.d.\ samplers.
\end{enumerate}
These are not merely aesthetic deficiencies.  Many biologically important quantities---kinetic rate
constants, mean first-passage times, committed fractions---are fundamentally temporal and cannot be
extracted from equilibrium samples alone.

\subsection{Static Structure Prediction Models}

\textbf{AlphaFold2}~\cite{jumper2021} achieves near-experimental accuracy in predicting the single
lowest-energy structure from an amino-acid sequence and a multiple-sequence alignment (MSA).
However, it produces a single structure per input; it was not designed to generate conformational
ensembles or dynamics.  MSA subsampling tricks can yield some alternative conformations, but the
variance is limited and not calibrated to MD statistics.

\textbf{ESMFold}~\cite{lin2023} extends the ESM-2 protein language model to structure prediction,
enabling fast single-structure prediction without MSA.  MDGen uses ESMFold (or any equivalent
structure predictor) only to generate the \emph{initial conditioning structure}; ESMFold itself
produces no ensemble.

\textbf{RFdiffusion}~\cite{watson2023} and \textbf{Chroma}~\cite{ingraham2023} apply diffusion
models to protein \emph{design}---generating novel backbone conformations that satisfy user-defined
functional constraints.  These models are not evaluated on MD ensemble fidelity metrics and are not
intended for dynamics prediction.

\subsection{Summary: What Was Missing Before MDGen}

\begin{table}[H]
  \centering
  \caption{Comparison of prior methods for protein ensemble characterisation.
           $\checkmark$ = capable; $\times$ = not capable; $\sim$ = partial.}
  \label{tab:prior_comparison}
  \renewcommand{\arraystretch}{1.3}
  \begin{tabularx}{\textwidth}{Xcccc}
    \toprule
    \textbf{Method} & \textbf{Equil.\ ensemble} & \textbf{Temporal coherence} & \textbf{Speed vs MD} & \textbf{Large proteins} \\
    \midrule
    MD simulation       & $\checkmark$ & $\checkmark$ & $1\times$ (baseline) & $\checkmark$ \\
    NMA                 & $\sim$       & $\times$     & $\sim 10^3\times$    & $\checkmark$ \\
    Boltzmann generators & $\checkmark$ & $\times$    & $\sim 10^3\times$    & $\times$ \\
    EigenFold / AlphaFlow & $\checkmark$ & $\times$  & $\sim 10^4\times$    & $\checkmark$ \\
    BioEmu              & $\checkmark$ & $\times$     & $\sim 10^4\times$    & $\checkmark$ \\
    \textbf{MDGen (this work)} & $\checkmark$ & $\checkmark$ & $\sim 10^4\times$ & $\checkmark$ \\
    \bottomrule
  \end{tabularx}
\end{table}

The defining novelty of MDGen is the first entry in Table~\ref{tab:prior_comparison} that achieves
both equilibrium ensemble quality \emph{and} temporal coherence at large-protein scale.


% ============================================================
\section{MDGen: Method}
\label{sec:mdgen}
% ============================================================

\subsection{Problem Formulation}

MDGen addresses the following generative modelling problem:

\begin{calloutbox}
  \textbf{Definition (Trajectory generation).}  Given a protein amino-acid sequence
  $\bs = (s_1, \ldots, s_N)$, generate a trajectory window
  $\bX = (\bx_1, \ldots, \bx_T) \in \bigl(\SE(3)^N \times \TT^{4N}\bigr)^T$ such that the joint
  distribution $p_\theta(\bX \mid \bs)$ approximates the distribution over $T$-frame windows drawn
  from equilibrium MD trajectories.
\end{calloutbox}

This formulation imposes two simultaneous requirements:
\begin{enumerate}[leftmargin=2em]
  \item \textbf{Equilibrium marginal correctness:} for any single time $t$, the marginal
    $p(\bx_t \mid \bs)$ matches the MD equilibrium distribution.
  \item \textbf{Temporal coherence:} the joint distribution over consecutive frames
    $p(\bx_t, \bx_{t+1} \mid \bs)$ reflects the physically correct Markov transition structure
    of MD dynamics.
\end{enumerate}

Training window size $T = 16$ frames, corresponding to approximately 16~ns of simulation at the
1~ns ATLAS frame rate.  At inference, trajectory windows can be chained by conditioning on the
last frame of the previous window (Section~\ref{sec:inpainting}).

\subsection{Forward Diffusion Process over Trajectory Space}
\label{sec:forward}

MDGen applies independent noise to each trajectory frame, keeping the forward process
factorised over the time axis.  For a clean trajectory $\bX^{(0)}$ at diffusion noise level
$s = 0$ and noisy trajectory $\bX^{(s)}$ at noise level $s \in [0, S]$:
\begin{equation}
  q(\bX^{(s)} \mid \bX^{(0)}) = \prod_{t=1}^{T} q(\bx_t^{(s)} \mid \bx_t^{(0)}),
  \label{eq:forward_trajectory}
\end{equation}
where each frame is noised independently.  Within a frame, diffusion is applied jointly but
factorised over geometric components:
\begin{equation}
  q(\bx_t^{(s)} \mid \bx_t^{(0)}) =
  q_{\RR^3}(\bt^{(s)} \mid \bt^{(0)}) \cdot
  q_{\SO(3)}(\bR^{(s)} \mid \bR^{(0)}) \cdot
  q_{\TT}(\bm\chi^{(s)} \mid \bm\chi^{(0)}),
  \label{eq:frame_forward}
\end{equation}
where:
\begin{itemize}[leftmargin=2em]
  \item $q_{\RR^3}$ is the standard Gaussian process on translations \eqref{eq:forward_marginal}.
  \item $q_{\SO(3)}$ is the IGSO3 process \eqref{eq:igso3}.
  \item $q_{\TT}$ is the wrapped Gaussian process \eqref{eq:wrapped_gaussian} applied independently
    to each torsion angle.
\end{itemize}

A critical observation: although frames are noised independently, the reverse process denoises
\emph{all $T$ frames jointly}.  Temporal correlations are thus \emph{not} encoded in the forward
process noise kernel; rather, they are entirely a property of the \emph{learned denoising
network}.  The network must simultaneously recover each frame's geometry \emph{and} learn the
temporal consistency constraints between frames.

\subsection{Neural Architecture}
\label{sec:architecture}

The denoising network $\beps_\theta(\bX^{(s)}, s, \bs)$ takes as input the noisy trajectory
$\bX^{(s)}$, the diffusion noise level $s$, and the protein sequence $\bs$, and outputs predicted
noise tensors for translations, rotations, and torsions at every residue and every frame.

The architecture consists of four main components:

\subsubsection{Sequence Encoder: ESM-2}

The protein sequence $\bs$ is embedded by the pre-trained protein language model
ESM-2~\cite{lin2023} into per-residue feature vectors:
\begin{equation}
  \bm{h}_i^{\text{seq}} = \text{ESM-2}(\bs)_i \in \RR^{d_{\text{seq}}}.
\end{equation}
ESM-2 is kept \emph{frozen} during MDGen training.  Its representations capture evolutionary and
physicochemical information about each residue in its sequence context, serving as a rich
conditioning signal.  The embedding dimension $d_{\text{seq}}$ is projected to the internal
model dimension $d$ via a learned linear layer.

\subsubsection{Spatial Attention: Invariant Point Attention (IPA)}

At each trajectory frame $t$, inter-residue information exchange is mediated by \textbf{Invariant
Point Attention (IPA)}~\cite{jumper2021}.  IPA augments standard multi-head attention with
\emph{3D query/key point pairs} expressed in each residue's local backbone frame.  The attention
score between residues $i$ and $j$ is:
\begin{equation}
  a_{ij} = \text{softmax}_j\!\left(\frac{\bm{q}_i^\top \bm{k}_j}{\sqrt{d}}
  - \frac{\gamma}{2} \sum_{h=1}^{H_{\text{pt}}}
  \|T_i \bm{p}_i^{(h)} - T_j \bm{p}_j^{(h)}\|^2\right),
  \label{eq:ipa}
\end{equation}
where:
\begin{itemize}[leftmargin=2em]
  \item $\bm{q}_i, \bm{k}_j \in \RR^d$ are the standard query/key vectors.
  \item $\bm{p}_i^{(h)} \in \RR^3$ are learned query 3D points in residue $i$'s local frame.
  \item $T_i = (\bR_i, \bt_i)$ maps local frame points to the global coordinate system.
  \item The term $\|T_i \bm{p}_i^{(h)} - T_j \bm{p}_j^{(h)}\|^2$ measures the Euclidean
    distance between point $h$ in frame $i$'s frame and the corresponding point in frame $j$'s
    frame, \emph{after transforming both to global coordinates}.
  \item $\gamma > 0$ is a learnable scalar that controls the relative weight of geometric distance
    vs.\ sequence-feature similarity.
\end{itemize}

The geometry term makes the attention \textbf{SE(3)-invariant}: because $\|T_i \bm{p} - T_j
\bm{q}\|$ is unchanged by any global rotation or translation applied to all frames simultaneously,
the attention scores are physically meaningful and independent of the arbitrary lab-frame
orientation.  Crucially, two residues that are distant in sequence but close in 3D space receive a
high attention weight, enabling the model to reason about spatial contacts.

IPA is applied independently at each trajectory frame $t$ (``spatial pass''), allowing all
$N$ residues to exchange information conditioned on their current noisy geometry.

\subsubsection{Temporal Attention: Axial Attention Across Time}
\label{sec:temporal_attention}

The key architectural innovation in MDGen is the introduction of \textbf{axial temporal attention}~\cite{ho2019axial}.
After the spatial IPA pass (which operates at fixed $t$, across residues), a temporal attention
pass operates at fixed residue $i$, across all $T$ time steps:
\begin{equation}
  \bm{h}_{i,t}^{\text{new}} = \sum_{t'=1}^{T} \alpha_{tt'}\,\bm{v}_{i,t'},
  \qquad
  \alpha_{tt'} = \text{softmax}_{t'}\!\!\left(\frac{\bm{q}_{i,t} \cdot \bm{k}_{i,t'}}{\sqrt{d}}\right),
  \label{eq:temporal_attn}
\end{equation}
where $\bm{q}_{i,t}, \bm{k}_{i,t'}, \bm{v}_{i,t'} \in \RR^d$ are the query, key, and value
vectors for residue $i$ at times $t$ and $t'$, respectively.

The temporal attention layer allows every time step to \emph{attend to every other time step}
within the trajectory window, learning which temporal relationships are informative for denoising.
This is precisely the mechanism through which MDGen learns that consecutive frames should be
physically similar: during training, the loss penalises temporally incoherent predictions, and the
temporal attention weights adapt to propagate continuity information along the time axis.

The spatial and temporal attention layers alternate in a stack of $L$ blocks:
\begin{equation}
  \text{Block}_\ell:\quad
  \underbrace{\text{IPA}_\ell(\cdot)}_{\text{spatial, per frame}} \;\to\;
  \underbrace{\text{TempAttn}_\ell(\cdot)}_{\text{temporal, per residue}} \;\to\;
  \underbrace{\text{FFN}_\ell(\cdot)}_{\text{feed-forward}},
\end{equation}
constituting an \emph{axial attention transformer} over the $N \times T$ grid of (residue, frame)
tokens.

\subsubsection{Score Head}

After $L$ alternating blocks, separate prediction heads convert the latent representation
$\bm{h}_{i,t}^{(L)}$ into predicted noise for each geometric component:
\begin{align}
  \hat\beps^{\text{trans}}_{i,t} &= \text{MLP}_{\text{trans}}(\bm{h}_{i,t}^{(L)}) \in \RR^3, \\
  \hat\beps^{\text{rot}}_{i,t} &= \text{MLP}_{\text{rot}}(\bm{h}_{i,t}^{(L)}) \in \mathfrak{so}(3), \\
  \hat\beps^{\text{tors}}_{i,t} &= \text{MLP}_{\text{tors}}(\bm{h}_{i,t}^{(L)}) \in \RR^4,
\end{align}
where the rotation score is expressed as a vector in the tangent space $\mathfrak{so}(3) \cong
\RR^3$ of the Lie algebra.  The total parameter count is approximately 100M.

\subsubsection{Architecture Summary}

\begin{algorithm}[H]
  \caption{MDGen Forward Pass}
  \begin{algorithmic}[1]
    \Require Noisy trajectory $\bX^{(s)} = \{(\bR_{i,t}^{(s)}, \bt_{i,t}^{(s)}, \bm\chi_{i,t}^{(s)})\}$,
             noise level $s$, sequence $\bs$
    \State $\bm{h}_{i}^{\text{seq}} \gets \text{ESM-2}(\bs)_i$ for all $i$ \Comment{frozen}
    \State $\bm{h}_{i,t} \gets \text{Linear}([\bm{h}_i^{\text{seq}}; \phi(s); \text{frame\_embed}(\bR_{i,t}^{(s)}, \bt_{i,t}^{(s)})])$
    \For{$\ell = 1, \ldots, L$}
      \For{$t = 1, \ldots, T$} \Comment{spatial pass, parallelised}
        \State $\bm{h}_{i,t} \gets \text{IPA}_\ell\bigl(\{\bm{h}_{j,t}\}_{j=1}^N, \{T_{j,t}\}_{j=1}^N\bigr)$ for all $i$
      \EndFor
      \For{$i = 1, \ldots, N$} \Comment{temporal pass, parallelised}
        \State $\bm{h}_{i,t} \gets \text{TempAttn}_\ell\bigl(\{\bm{h}_{i,t'}\}_{t'=1}^T\bigr)$ for all $t$
      \EndFor
      \State $\bm{h}_{i,t} \gets \text{FFN}_\ell(\bm{h}_{i,t})$ for all $i,t$
    \EndFor
    \State \Return $\hat\beps^{\text{trans}}, \hat\beps^{\text{rot}}, \hat\beps^{\text{tors}}$ \Comment{via score heads}
  \end{algorithmic}
\end{algorithm}

\subsection{Training Objective}

The training loss is a weighted sum of denoising score matching losses for each geometric component,
plus a structural violation penalty:
\begin{equation}
  \mathcal{L} = \EE_{s,\bX^{(0)},\bX^{(s)}}\!\Bigl[
    \lambda_{\text{trans}}\,\mathcal{L}_{\RR^3}
    + \lambda_{\text{rot}}\,\mathcal{L}_{\SO(3)}
    + \lambda_{\text{tors}}\,\mathcal{L}_{\TT}
    + \lambda_{\text{viol}}\,\mathcal{L}_{\text{viol}}
  \Bigr],
  \label{eq:total_loss}
\end{equation}
where each component is:

\textbf{Translation loss:}
\begin{equation}
  \mathcal{L}_{\RR^3} = \frac{1}{NT}\sum_{i,t} \|\hat\beps_{i,t}^{\text{trans}} - \beps_{i,t}^{\text{trans}}\|^2.
  \label{eq:trans_loss}
\end{equation}

\textbf{Rotation loss} (Riemannian score matching on $\SO(3)$):
\begin{equation}
  \mathcal{L}_{\SO(3)} = \frac{1}{NT}\sum_{i,t}
  \|\hat\beps_{i,t}^{\text{rot}} - \beps_{i,t}^{\text{rot}}\|^2,
  \label{eq:rot_loss}
\end{equation}
where $\beps_{i,t}^{\text{rot}}$ is the score in the Lie algebra at the noisy rotation, computed
analytically from the IGSO3 forward process.

\textbf{Torsion loss} (score matching on $\TT$):
\begin{equation}
  \mathcal{L}_{\TT} = \frac{1}{4NT}\sum_{i,t,k}
  \bigl(\hat\epsilon_{i,t,k}^{\text{tors}} - \epsilon_{i,t,k}^{\text{tors}}\bigr)^2,
  \label{eq:tors_loss}
\end{equation}
where the sum is over the four torsion angles $k \in \{1,2,3,4\}$ and the score
$\epsilon_{i,t,k}^{\text{tors}}$ is computed from the wrapped Gaussian process.

\textbf{Violation loss:}
\begin{equation}
  \mathcal{L}_{\text{viol}} = \mathcal{L}_{\text{clash}} + \mathcal{L}_{\text{bond}},
  \label{eq:viol_loss}
\end{equation}
comprising a pairwise steric clash penalty and a backbone bond-length/angle violation penalty.
This term is small ($\lambda_{\text{viol}} \approx 0.01$) but regularises the model toward
physically valid geometries.  The approximate loss weights used in the original implementation are
$\lambda_{\text{trans}} = 1.0$, $\lambda_{\text{rot}} \approx 1.0$, $\lambda_{\text{tors}}
\approx 0.5$, $\lambda_{\text{viol}} \approx 0.01$.

All $N$ residues and all $T = 16$ frames contribute to each gradient update, providing dense
training signal per data sample.  The diffusion noise level $s$ is sampled uniformly from
$[0, S]$ with $S = 500$ at each training step.

\textbf{Optimisation.}  MDGen is trained with the AdamW
optimiser~\cite{loshchilov2019}, learning rate $10^{-4}$ with cosine annealing, batch
size $\approx 64$ trajectory windows, on $8 \times$ NVIDIA A100 GPUs.

\subsection{Sampling / Inference}

\textbf{Unconditional generation.}  Starting from pure noise $\bX^{(S)}$ (translations
$\sim\mathcal{N}(\bm{0}, \bm{I})$, rotations $\sim \text{Uniform}(\SO(3))$, torsions
$\sim \text{Uniform}(\TT^4)$), the reverse diffusion process iterates from $s = S$ down to $s = 0$:
\begin{equation}
  \bX^{(s-1)} \gets \text{Reverse}(\bX^{(s)},\; \beps_\theta(\bX^{(s)}, s, \bs)),
\end{equation}
applying the reverse update separately for each geometric component (Gaussian reverse for
translations, IGSO3 reverse for rotations, wrapped Gaussian reverse for torsions).

Two sampling modes are supported:
\begin{itemize}[leftmargin=2em]
  \item \textbf{DDPM} (500 steps): stochastic reverse process; highest sample quality.
  \item \textbf{DDIM} (50 steps): deterministic ODE \eqref{eq:ddim} on each component;
    $10\times$ faster than DDPM with minimal quality degradation.
\end{itemize}
Both modes produce trajectories approximately $10{,}000\times$ faster than equivalent
explicit-solvent MD.

\subsection{Trajectory Extension via Diffusion Inpainting}
\label{sec:inpainting}

MDGen can generate trajectories conditioned on a known conformation $\bx_1^{(0)}$
(e.g.\ a crystal structure or AlphaFold2 prediction) using a technique inspired by image
inpainting~\cite{lugmayr2022}.  During the reverse diffusion process, the first-frame
variable is periodically replaced with its correctly-noised counterpart:
\begin{equation}
  \bx_1^{(s)} \leftarrow \sqrt{\bar\alpha_s}\,\bx_1^{(0)} + \sqrt{1-\bar\alpha_s}\,\beps,
  \label{eq:inpainting}
\end{equation}
while the remaining frames $\bx_2^{(s)}, \ldots, \bx_T^{(s)}$ evolve freely.  The temporal
attention mechanism propagates structural information from the conditioned frame to all others,
producing a physically consistent continuation trajectory.

\textbf{Interpolation mode.}  By conditioning on \emph{both} the first frame $\bx_1^{(0)}$ and
the last frame $\bx_T^{(0)}$, MDGen generates a plausible transition pathway between two known
protein conformations (e.g.\ open and closed states of an enzyme).  This capability for
\emph{transition path sampling} without enhanced sampling MD is a direct consequence of the
temporal attention architecture and has no equivalent in any prior independent-frame sampler.


% ============================================================
\section{Training Data: The ATLAS Database}
\label{sec:data}
% ============================================================

MDGen is trained on the \textbf{ATLAS} (Atlas of Trajectory Landscapes and Associated
Simulations) database~\cite{vandermeersche2024}, the largest open repository of protein MD
trajectories assembled specifically to support machine learning on protein dynamics.

\textbf{Simulation protocol.}
Each protein was simulated using explicit TIP3P water~\cite{jorgensen1983} with the AMBER
ff14SB force field~\cite{maier2015} under NPT conditions (constant number of particles, pressure,
and temperature) at 300~K and 1~atm.  Three independent 100~ns replicas were run per protein,
yielding 300~ns of total simulation time per protein.  Frames were saved at 1~ns intervals,
producing $\sim$300 frames per protein.  The simulation time step was 2~fs with SHAKE constraints
on hydrogen bonds.

\textbf{Dataset statistics.}
\begin{itemize}[leftmargin=2em]
  \item $\sim$1,400 diverse proteins from the PDB, spanning all major fold classes ($\alpha$,
    $\beta$, $\alpha/\beta$, and intrinsically disordered regions).
  \item Proteins range from $\sim$50 to $\sim$500 residues.
  \item Aggregate simulation time: $\sim$1.2~$\mu$s across all proteins and replicas.
\end{itemize}

\textbf{Train/test split.}
$\sim$1,265 proteins were used for training; $\sim$135 held out for test evaluation.  The test
set was constructed to assess generalisation to unseen proteins, including a subset of
well-characterised ``fast-folding'' miniproteins (Trp-cage, BBA, villin headpiece, WW domain,
NTL9, Chignolin) from long reference simulations by Lindorff-Larsen et al.~\cite{lindorff2011}
and D.E. Shaw Research that are not in ATLAS.

\textbf{Trajectory windowing.}  During training, a random contiguous window of $T = 16$
consecutive frames is extracted from each replica trajectory for each minibatch example.  Overlapping
windows are permitted.  This augmentation strategy exposes the model to the full range of dynamical
timescales present in the data.


% ============================================================
\section{Evaluation Protocol and Results}
\label{sec:results}
% ============================================================

\subsection{Evaluation Metrics}

MDGen is evaluated on five categories of metrics, carefully chosen to assess both equilibrium and
kinetic fidelity:

\subsubsection{Root Mean Square Fluctuation (RMSF)}

RMSF quantifies per-residue positional flexibility:
\begin{equation}
  \text{RMSF}_i = \sqrt{\frac{1}{T}\sum_{t=1}^{T}\|\br_i(t) - \langle\br_i\rangle\|^2},
  \label{eq:rmsf}
\end{equation}
where $\br_i(t)$ is the $\text{C}_\alpha$ position of residue $i$ at frame $t$ and $\langle\br_i\rangle
= \frac{1}{T}\sum_t \br_i(t)$ is its temporal mean position.  RMSF is computed per-residue and
reported as a profile over the sequence.

The primary summary metric is the \textbf{Pearson correlation coefficient} $r$ between the MDGen
RMSF profile and the reference MD RMSF profile, averaged over all test proteins.  A high $r$
indicates that MDGen correctly identifies which residues are flexible (loops, termini) and which
are rigid (helices, $\beta$-sheets).

\subsubsection{Ramachandran Distribution: Jensen--Shannon Divergence}

The Ramachandran plot captures the joint distribution of backbone dihedral angles $(\phi, \psi)$
for all residues across the ensemble.  MDGen's generated $(\phi, \psi)$ distribution is compared
to the reference MD distribution using the \textbf{Jensen--Shannon divergence (JSD)}:
\begin{equation}
  \text{JSD}(P \| Q) = \frac{1}{2} \text{KL}(P \| M) + \frac{1}{2} \text{KL}(Q \| M),
  \quad M = \frac{P + Q}{2},
  \label{eq:jsd}
\end{equation}
where $P$ is the reference MD distribution and $Q$ is the generated distribution, both discretised
over a $36 \times 36$ grid of $(\phi, \psi)$ bins.  Lower JSD indicates better agreement.

\subsubsection{Frame-to-Frame RMSD Distribution}

For each pair of consecutive frames $(t, t+1)$, the $\text{C}_\alpha$ RMSD between frames is
computed:
\begin{equation}
  \Delta_t = \text{RMSD}\bigl(\bm{r}^{(t+1)}_{\text{C}_\alpha}, \bm{r}^{(t)}_{\text{C}_\alpha}\bigr).
  \label{eq:frame_rmsd}
\end{equation}
The distribution of $\Delta_t$ over the test set quantifies whether MDGen respects the
\emph{physical step size} of MD dynamics.  In real MD, $\Delta_t$ is tightly distributed around
$\sim$0.1--0.3~\AA; an independent sampler produces $\Delta_t$ values many times larger.

\subsubsection{Temporal Autocorrelation Function (ACF)}

The autocorrelation function measures the temporal persistence of structural fluctuations:
\begin{equation}
  C(\tau) = \frac{\langle \delta\bm{r}(t) \cdot \delta\bm{r}(t+\tau) \rangle}
                 {\langle |\delta\bm{r}(t)|^2 \rangle},
  \label{eq:acf}
\end{equation}
where $\delta\bm{r}(t) = \bm{r}(t) - \langle\bm{r}\rangle$ is the displacement from the mean
structure.  Physical trajectories exhibit a smooth decay from $C(0) = 1$ to a baseline,
reflecting the timescales of structural relaxation.  An independent sampler has $C(\tau) = 0$ for
all $\tau \geq 1$ by definition.

\subsubsection{Side-Chain Torsion Distributions}

The $\chi_1, \chi_2$ torsion angle distributions are compared per-residue type using JSD, assessing
whether the generated side-chain conformations match the rotamer populations observed in MD.

\subsection{Quantitative Results}

\begin{table}[H]
  \centering
  \caption{Quantitative evaluation of MDGen and baselines on held-out ATLAS test proteins.
           Values are approximate; see Jing et al.~\cite{jing2024mdgen} for exact numbers
           and confidence intervals.}
  \label{tab:results}
  \renewcommand{\arraystretch}{1.3}
  \begin{tabularx}{\textwidth}{Xccccc}
    \toprule
    \textbf{Method} & \textbf{RMSF $r$ $\uparrow$} & \textbf{Ramach.\ JSD $\downarrow$}
      & \textbf{Frame $\Delta$ RMSD} & \textbf{ACF match} & \textbf{Speed vs MD} \\
    \midrule
    Reference MD                  & 1.00 & 0.00 & ground truth     & ground truth  & $1\times$ \\
    \textbf{MDGen} (DDPM)         & \textbf{$\approx$0.87} & \textbf{$\approx$0.08} & \checkmark matches & \checkmark matches & $\approx 10^4\times$ \\
    MDGen (DDIM)                  & $\approx$0.85 & $\approx$0.09 & \checkmark matches & \checkmark matches & $\approx 10^5\times$ \\
    Indep.\ frame diffusion       & $\approx$0.71 & $\approx$0.11 & $\times$ random   & $\times$ zero lag  & $\approx 10^4\times$ \\
    Normal Mode Analysis          & $\approx$0.64 & $\approx$0.19 & harmonic only     & harmonic only      & $\approx 10^3\times$ \\
    Boltzmann generators          & $\approx$0.68 & $\approx$0.14 & $\times$ i.i.d.   & $\times$ zero lag  & $\approx 10^3\times$ \\
    \bottomrule
  \end{tabularx}
\end{table}

\textbf{RMSF.}  MDGen achieves Pearson $r \approx 0.87$ on per-residue RMSF, substantially
outperforming both NMA ($r \approx 0.64$) and independent-frame diffusion ($r \approx 0.71$).
The improvement over independent-frame diffusion indicates that temporal context improves
equilibrium ensemble quality as well as kinetics---consistent with the interpretation that
temporal attention acts as an informative prior constraining the physically accessible region of
conformation space.

\textbf{Ramachandran distributions.}  MDGen achieves JSD $\approx 0.08$ compared to the reference
MD Ramachandran distribution, versus NMA's JSD $\approx 0.19$.  NMA's high JSD reflects the
harmonic approximation's failure to reproduce anharmonic populations in loop regions and
$\beta$-turns.

\textbf{Temporal metrics.}  Only MDGen reproduces correct temporal autocorrelation functions and
frame-to-frame step-size distributions.  Independent-frame methods score zero on ACF match by
construction.

\subsection{Generalisation to Out-of-Distribution Fast-Folding Proteins}

To assess generalisation beyond the ATLAS training distribution, MDGen is evaluated on
miniproteins from long reference simulations by Lindorff-Larsen et al.~\cite{lindorff2011} that
were not in the ATLAS training set.  These include Trp-cage (20 residues), BBA motif (28
residues), villin headpiece (35 residues), WW domain (35 residues), NTL9 (39 residues), and
Chignolin (10 residues).

MDGen correctly identifies the qualitative flexibility profile of these proteins: flexible C-terminal
helices, rigid hydrophobic cores, and loop regions with elevated RMSF.  The RMSF profiles closely
track the reference simulations, demonstrating that the model generalises across protein fold space
beyond the globular, single-chain proteins that dominate ATLAS.


% ============================================================
\section{Applications}
\label{sec:applications}
% ============================================================

\subsection{Cryptic Pocket Discovery in Drug Design}

Approximately 50\% of protein drug targets undergo conformational changes critical for inhibitor
binding.  A substantial fraction of binding sites are \emph{cryptic}---they do not exist in the
static crystal structure but appear transiently during conformational dynamics~\cite{oleinikovas2016}.
Static docking against a single experimental structure misses these pockets entirely, contributing
to high false-negative rates in virtual screening.

MDGen enables a practical workflow for cryptic pocket discovery:
\begin{enumerate}[leftmargin=2em]
  \item Generate $10^2$--$10^3$ backbone trajectory frames from the protein sequence using MDGen
    (seconds to minutes on a single GPU).
  \item Run pocket detection software (FPOCKET~\cite{le2009}, SiteMap) on each frame.
  \item Identify pockets that appear transiently across the ensemble.
  \item Prioritise pockets for downstream docking and lead optimisation.
\end{enumerate}

This approach mirrors ensemble docking~\cite{totrov2008}, but replaces expensive multi-$\mu$s MD
with MDGen-generated ensembles that are three to four orders of magnitude faster to produce.
MDGen's correct temporal autocorrelation further ensures that the generated ensemble samples
conformational transitions consistent with the protein's natural dynamics rather than an arbitrary
collection of unrelated structures.

\subsection{Intrinsically Disordered Proteins (IDPs) and Disease}

Intrinsically disordered proteins (IDPs)~\cite{uversky2019} lack a stable folded structure and
exist as highly dynamic conformational ensembles.  They are biologically ubiquitous---involved in
transcription regulation, signalling, and cell division---and medically critical, as their
disordered regions frequently drive aberrant condensate formation and aggregation in
neurodegenerative diseases.

Key disease-relevant IDPs include:
\begin{itemize}[leftmargin=2em]
  \item \textbf{$\alpha$-Synuclein (Parkinson's disease):} an $\sim$140-residue IDP whose
    conformational ensemble determines the nucleation and growth pathways of toxic amyloid fibrils.
  \item \textbf{Tau (Alzheimer's disease):} the highly flexible repeat domain of tau adopts diverse
    conformations; specific conformational states initiate neurofibrillary tangle formation.
  \item \textbf{p53 (cancer):} the most frequently mutated tumour suppressor gene; its disordered
    N- and C-terminal regions mediate transactivation and regulate DNA binding through
    conformational dynamics.
\end{itemize}

Static predictors such as AlphaFold2 report low-confidence (pLDDT $< 50$) predictions for
disordered regions, which is a diagnostic flag rather than a structural answer.  MDGen generates
backbone conformational ensembles for these regions by leveraging the sequence-conditioned
generative model, producing the distribution of conformations rather than a misleading single
coordinate set.

\subsection{Ensemble-Aware Protein Engineering}

Beyond drug discovery, MDGen's trajectory generation capability enables:
\begin{itemize}[leftmargin=2em]
  \item \textbf{Mutational scanning:} rapidly generate conformational ensembles for point mutants
    to predict effects on dynamics, stability, and allostery without re-running MD for each variant.
  \item \textbf{Allosteric pathway mapping:} condition on two allosteric states (active/inactive)
    and use interpolation mode (Section~\ref{sec:inpainting}) to identify residues involved in
    signal transmission.
  \item \textbf{Antibody engineering:} characterise the CDR loop flexibility that governs antigen
    recognition and affinity.
\end{itemize}


% ============================================================
\section{Limitations}
\label{sec:limitations}
% ============================================================

While MDGen represents a significant advance in generative protein dynamics modelling, several
important limitations must be acknowledged:

\begin{enumerate}[leftmargin=2em]
  \item \textbf{Fixed timescale window.}  The training window of $T = 16$ frames corresponds to
    approximately 16~ns at the ATLAS 1~ns frame rate.  Millisecond-scale processes such as full
    folding/unfolding transitions, large-scale domain translocations, or cryptic conformational
    changes that occur on $\mu$s--ms timescales are outside the scope of a single inference
    call.  Window chaining via inpainting can extend coverage, but whether kinetic accuracy is
    preserved across chains requires further investigation.

  \item \textbf{Training distribution bias.}  ATLAS predominantly contains single-chain, globular,
    soluble proteins of moderate size (50--500 residues).  Membrane proteins, large multi-domain
    assemblies, oligomeric complexes, and heavily post-translationally modified proteins are
    under-represented or absent.  Generalisation to these classes is not guaranteed.

  \item \textbf{Implicit solvent statistics.}  MDGen learns the joint statistics of protein
    coordinates from explicit-solvent MD but does not model the solvent molecules explicitly.
    Hydration shell dynamics, specific ion binding, and osmolyte effects are implicitly captured
    only insofar as they manifest in backbone motions in the training trajectories.

  \item \textbf{Force-field dependency.}  Generated trajectories emulate the statistical properties
    of AMBER ff14SB~\cite{maier2015}.  Systematic errors in this force field (e.g.\ known
    overstabilisation of $\alpha$-helices in some variants) will be inherited by MDGen.  The model
    is not an independent physics engine; it is a data-driven emulator of a specific force field
    on a specific dataset.

  \item \textbf{Backbone-only representation.}  MDGen generates backbone frames and side-chain
    torsion angles, but does not produce full all-atom coordinates directly.  Side-chain
    coordinates must be reconstructed post-hoc using tools such as Rosetta Relax~\cite{conway2014}
    or DLPacker~\cite{xiong2022}.  This limits direct use of generated trajectories in applications
    requiring atom-level precision (e.g.\ binding free energy calculations by alchemical
    perturbation).

  \item \textbf{Rare barrier crossings.}  Conformational transitions that require crossing high
    energy barriers (rare events on $\mu$s+ timescales) are statistically improbable within a
    16-frame window.  MDGen is therefore primarily a tool for characterising fast, accessible
    dynamics near the dominant energy basin(s)---not for studying rare-event kinetics.
\end{enumerate}


% ============================================================
\section{Future Directions}
\label{sec:future}
% ============================================================

MDGen establishes a proof-of-concept for generative molecular dynamics, opening a rich set of
future research directions:

\paragraph{Full-atom trajectory generation.}
Extending MDGen to jointly generate backbone frames and all-atom side-chain coordinates as a
unified diffusion model would enable direct computation of binding free energies, hydrogen-bond
networks, and pharmacophoric features without post-hoc side-chain reconstruction.

\paragraph{Longer timescales via enhanced training data.}
The 300~ns ATLAS trajectories limit MDGen to ns-scale dynamics.  Training on $\mu$s-scale datasets
from D.E. Shaw Research's Anton supercomputer~\cite{shaw2010} or on enhanced-sampling
trajectories (metadynamics, REST2) would extend the accessible timescale and enable modelling of
rarer conformational transitions.

\paragraph{Protein--ligand and protein--protein dynamics.}
Extending the SE(3) frame representation to multi-chain and receptor--ligand systems would enable
ensemble-aware docking, entropic contributions to binding affinity ($\Delta G = \Delta H - T\Delta
S$), and dynamics at protein--protein interfaces.

\paragraph{Conditioning on experimental observables.}
Incorporating experimental restraints from NMR chemical shifts, SAXS scattering profiles, or
cryo-EM density maps as soft conditioning signals during inference would enable
\emph{physics-guided ensemble refinement}---generating ensembles consistent with both the learned
dynamics prior and experimental measurements.

\paragraph{Retraining on alternative force fields.}
The MDGen framework is agnostic to the training data source.  Retraining on trajectories generated
with polarisable force fields, QM/MM potentials, or coarse-grained models would yield surrogates
tuned to different levels of physical accuracy.

\paragraph{Beyond proteins.}
The SE(3) trajectory diffusion framework is not protein-specific.  RNA, DNA, carbohydrates, and
lipid membranes all undergo important conformational dynamics accessible to MD; the same
architectural principles could be adapted with appropriate geometric representations for these
classes.

\paragraph{Latent dynamics and interpretability.}
Analysing the temporal attention weights of trained MDGen models may reveal which residues and
timescales dominate the information flow during trajectory generation, providing interpretable
insights into allosteric mechanisms and kinetic pathways.


% ============================================================
\section{Conclusion}
% ============================================================

MDGen represents a conceptual advance in the emerging field of generative molecular modelling.
By framing the problem as \emph{trajectory generation} rather than independent equilibrium
sampling, MDGen is the first model to simultaneously achieve:
\begin{itemize}[leftmargin=2em]
  \item High per-residue RMSF correlation with reference MD ($r \approx 0.87$, versus NMA $r
    \approx 0.64$).
  \item Physically correct Ramachandran backbone dihedral distributions (JSD $\approx 0.08$).
  \item Correct temporal autocorrelation functions and frame-to-frame step-size distributions---
    metrics that no prior independent-frame sampler can reproduce.
  \item A $\sim 10{,}000\times$ speed-up over equivalent explicit-solvent MD simulation.
\end{itemize}

The architectural innovations enabling this---the combination of SE(3)-equivariant Invariant Point
Attention (spatial) with temporal axial attention (kinetic)---provide a template for incorporating
physical symmetries and temporal structure into generative models for other molecular and
biological systems.  The framework is trained on the ATLAS database of 1,400 diverse 300~ns MD
trajectories and generalises to out-of-distribution fast-folding proteins not present in training.

MDGen's central intellectual contribution is demonstrating that MD trajectories should be treated
as the \emph{primary generative target}---not a source of independent equilibrium samples.  This
reframing enables models to learn both the thermodynamics (what conformations are accessible) and
the dynamics (how fast they interconvert) of protein motion from data.  It is a template for how
generative AI can emulate, rather than merely approximate, the complex physical processes that
govern molecular biology.


% ============================================================
% References
% ============================================================
\newpage
\begin{thebibliography}{99}

\bibitem{jing2024mdgen}
Jing, B., St\"{a}rk, H., Jaakkola, T., \& Berger, B. (2024).
Generative Molecular Dynamics.
\textit{arXiv:2409.17808}. Advances in Neural Information Processing Systems (NeurIPS 2024).
Code: \url{https://github.com/bjing2016/mdgen}

\bibitem{vandermeersche2024}
Vander Meersche, Y., Cretin, G., de Brevern, A.G., Gelly, J.C., \& Postic, G. (2024).
ATLAS: Protein Flexibility Description from Atomistic Molecular Dynamics Simulations.
\textit{Nucleic Acids Research}, 52(D1), D384--D392.

\bibitem{jumper2021}
Jumper, J., Evans, R., Pritzel, A., Green, T., Figurnov, M., Ronneberger, O., \ldots \& Hassabis, D. (2021).
Highly accurate protein structure prediction with AlphaFold.
\textit{Nature}, 596(7873), 583--589.

\bibitem{ho2020}
Ho, J., Jain, A., \& Abbeel, P. (2020).
Denoising diffusion probabilistic models.
\textit{Advances in Neural Information Processing Systems}, 33, 6840--6851.

\bibitem{song2021}
Song, Y., Sohl-Dickstein, J., Kingma, D.P., Kumar, A., Ermon, S., \& Poole, B. (2021).
Score-based generative modeling through stochastic differential equations.
\textit{International Conference on Learning Representations (ICLR 2021)}.

\bibitem{song2021ddim}
Song, J., Meng, C., \& Ermon, S. (2021).
Denoising diffusion implicit models.
\textit{International Conference on Learning Representations (ICLR 2021)}.

\bibitem{lin2023}
Lin, Z., Akin, H., Rao, R., Hie, B., Zhu, Z., Lu, W., \ldots \& Rives, A. (2023).
Evolutionary-scale prediction of atomic-level protein structure with a language model.
\textit{Science}, 379(6637), 1123--1130.

\bibitem{watson2023}
Watson, J.L., Juergens, D., Bennett, N.R., Trippe, B.L., Yim, J., Eisenach, H.E., \ldots \& Baker, D. (2023).
De novo design of protein structure and function with RFdiffusion.
\textit{Nature}, 620(7976), 1089--1100.

\bibitem{adcock2006}
Adcock, S.A., \& McCammon, J.A. (2006).
Molecular dynamics: survey of methods for simulating the activity of proteins.
\textit{Chemical Reviews}, 106(5), 1589--1615.

\bibitem{shaw2010}
Shaw, D.E., Maragakis, P., Lindorff-Larsen, K., Piana, S., Dror, R.O., Eastwood, M.P., \ldots \& Wriggers, W. (2010).
Atomic-level characterization of the structural dynamics of proteins.
\textit{Science}, 330(6002), 341--346.

\bibitem{lindorff2011}
Lindorff-Larsen, K., Piana, S., Dror, R.O., \& Shaw, D.E. (2011).
How fast-folding proteins fold.
\textit{Science}, 334(6055), 517--520.

\bibitem{debortoli2022}
De Bortoli, V., Mathieu, E., Hutchinson, M., Thornton, J., Teh, Y.W., \& Doucet, A. (2022).
Riemannian score-based generative modelling.
\textit{Advances in Neural Information Processing Systems}, 35, 2406--2422.

\bibitem{leach2022}
Leach, A., Schmon, S.M., Degiacomi, M.T., \& Willcocks, C.G. (2022).
Denoising diffusion probabilistic models on SO(3) for rotational alignment.
\textit{ICLR 2022 Workshop on Geometrical and Topological Representation Learning}.

\bibitem{yim2023framediff}
Yim, J., Trippe, B.L., De Bortoli, V., Mathieu, E., Doucet, A., Barzilay, R., \& Jaakkola, T. (2023).
SE(3) diffusion model with application to protein backbone generation.
\textit{International Conference on Machine Learning (ICML 2023)}.

\bibitem{jing2023eigenfold}
Jing, B., Corso, G., Chang, J., Barzilay, R., \& Jaakkola, T. (2023).
EigenFold: Generative protein structure prediction with diffusion models.
\textit{ICLR 2023 Workshop on Machine Learning for Drug Discovery}.

\bibitem{jing2024alphaflow}
Jing, B., Berger, B., \& Jaakkola, T. (2024).
AlphaFold meets flow matching for generating protein ensembles.
\textit{International Conference on Machine Learning (ICML 2024)}.

\bibitem{lewis2025bioemu}
Lewis, S.M., Hempel, T., Jim\'{e}nez-Luna, J., Gastegger, M., Xie, Y., Foong, A.Y.K., \ldots \& No\'{e}, F. (2025).
Scalable emulation of protein equilibrium ensembles with generative deep learning.
\textit{Science}, 388(6748), 1--13.

\bibitem{noe2019}
No\'{e}, F., Olsson, S., K\"{o}hler, J., \& Wu, H. (2019).
Boltzmann generators: sampling equilibrium states of many-body systems with deep learning.
\textit{Science}, 365(6457), eaaw1147.

\bibitem{tirion1996}
Tirion, M.M. (1996).
Large amplitude elastic motions in proteins from a single-parameter, atomic analysis.
\textit{Physical Review Letters}, 77(9), 1905.

\bibitem{brooks1983}
Brooks, B., \& Karplus, M. (1983).
Harmonic dynamics of proteins: normal modes and fluctuations in bovine pancreatic trypsin inhibitor.
\textit{Proceedings of the National Academy of Sciences}, 80(21), 6571--6575.

\bibitem{maier2015}
Maier, J.A., Martinez, C., Kasavajhala, K., Wickstrom, L., Hauser, K.E., \& Simmerling, C. (2015).
ff14SB: improving the accuracy of protein side chain and backbone parameters from ff99SB.
\textit{Journal of Chemical Theory and Computation}, 11(8), 3696--3713.

\bibitem{case2021}
Case, D.A., Aktulga, H.M., Betz, K., Ben-Shalom, I.Y., Berryman, J.T., Borden, S.R., \ldots \& Kollman, P.A. (2021).
Amber 2021. University of California, San Francisco.

\bibitem{abraham2015}
Abraham, M.J., Murtola, T., Schulz, R., P\'{a}ll, S., Smith, J.C., Hess, B., \& Lindahl, E. (2015).
GROMACS: High performance molecular simulations through multi-level parallelism from laptops to supercomputers.
\textit{SoftwareX}, 1, 19--25.

\bibitem{frauenfelder1991}
Frauenfelder, H., Sligar, S.G., \& Wolynes, P.G. (1991).
The energy landscapes and motions of proteins.
\textit{Science}, 254(5038), 1598--1603.

\bibitem{sugita1999}
Sugita, Y., \& Okamoto, Y. (1999).
Replica-exchange molecular dynamics method for protein folding.
\textit{Chemical Physics Letters}, 314(1--2), 141--151.

\bibitem{laio2002}
Laio, A., \& Parrinello, M. (2002).
Escaping free-energy minima.
\textit{Proceedings of the National Academy of Sciences}, 99(20), 12562--12566.

\bibitem{hamelberg2004}
Hamelberg, D., Mongan, J., \& McCammon, J.A. (2004).
Accelerated molecular dynamics: a promising and efficient simulation method for biomolecules.
\textit{Journal of Chemical Physics}, 120(24), 11919--11929.

\bibitem{fuchs2020}
Fuchs, F.B., Worrall, D.E., Fischer, V., \& Welling, M. (2020).
SE(3)-transformers: 3D roto-translation equivariant attention networks.
\textit{Advances in Neural Information Processing Systems}, 33, 1970--1981.

\bibitem{satorras2021}
Satorras, V.G., Hoogeboom, E., \& Welling, M. (2021).
E(n) equivariant graph neural networks.
\textit{International Conference on Machine Learning (ICML 2021)}.

\bibitem{jorgensen1983}
Jorgensen, W.L., Chandrasekhar, J., Madura, J.D., Impey, R.W., \& Klein, M.L. (1983).
Comparison of simple potential functions for simulating liquid water.
\textit{Journal of Chemical Physics}, 79(2), 926--935.

\bibitem{vincent2011}
Vincent, P. (2011).
A connection between score matching and denoising autoencoders.
\textit{Neural Computation}, 23(7), 1661--1674.

\bibitem{anderson1982}
Anderson, B.D.O. (1982).
Reverse-time diffusion equation models.
\textit{Stochastic Processes and their Applications}, 12(3), 313--326.

\bibitem{loshchilov2019}
Loshchilov, I., \& Hutter, F. (2019).
Decoupled weight decay regularization.
\textit{International Conference on Learning Representations (ICLR 2019)}.

\bibitem{lugmayr2022}
Lugmayr, A., Danelljan, M., Romero, A., Yu, F., Timofte, R., \& Van Gool, L. (2022).
Repaint: Inpainting using denoising diffusion probabilistic models.
\textit{IEEE Conference on Computer Vision and Pattern Recognition (CVPR 2022)}.

\bibitem{jing2022torsional}
Jing, B., Corso, G., Chang, J., Barzilay, R., \& Jaakkola, T. (2022).
Torsional diffusion for molecular conformer generation.
\textit{Advances in Neural Information Processing Systems}, 35, 24240--24253.

\bibitem{ho2019axial}
Ho, J., Kalchbrenner, N., Weissenborn, D., \& Salimans, T. (2019).
Axial attention in multidimensional transformers.
\textit{arXiv:1912.12180}.

\bibitem{oleinikovas2016}
Oleinikovas, V., Saladino, G., Cossins, B.P., \& Gervasio, F.L. (2016).
Understanding cryptic pocket formation in protein targets by enhanced sampling simulations.
\textit{Journal of the American Chemical Society}, 138(43), 14257--14263.

\bibitem{totrov2008}
Totrov, M., \& Abagyan, R. (2008).
Flexible ligand docking to multiple receptor conformations: a practical alternative.
\textit{Current Opinion in Structural Biology}, 18(2), 178--184.

\bibitem{le2009}
Le Guilloux, V., Schmidtke, P., \& Tuffery, P. (2009).
Fpocket: an open source platform for ligand pocket detection.
\textit{BMC Bioinformatics}, 10(1), 168.

\bibitem{uversky2019}
Uversky, V.N. (2019).
Intrinsically disordered proteins and their ``mysterious'' (meta)physics.
\textit{Frontiers in Physics}, 7, 10.

\bibitem{conway2014}
Conway, P., Tyka, M.D., DiMaio, F., Konerding, D.E., \& Baker, D. (2014).
Relaxation of backbone bond geometry improves protein energy landscape modeling.
\textit{Protein Science}, 23(1), 47--55.

\bibitem{xiong2022}
Xiong, P., Zhang, C., Zheng, W., \& Zhang, Y. (2022).
Protein side-chain conformation prediction by neural network with geometrically accurate atomic features.
\textit{Frontiers in Molecular Biosciences}, 8, 781038.

\bibitem{ingraham2023}
Ingraham, J.B., Baranov, M., Costello, Z., Barber, K.W., Wang, W., Ismail, A., \ldots \& Grigoryan, G. (2023).
Illuminating protein space with a programmable generative model.
\textit{Nature}, 623(7989), 1070--1078.

\bibitem{boehr2009}
Boehr, D.D., Nussinov, R., \& Wright, P.E. (2009).
The role of dynamic conformational ensembles in biomolecular recognition.
\textit{Nature Chemical Biology}, 5(11), 789--796.

\end{thebibliography}

\end{document}
